\section{Semantic memory and the status of a proposition}\label{sec:C14:semantic}
Semantic memory attempts to represent reusable information without carrying the complete episode each time. Its advantage is efficient access to facts, preferences, and abstractions. Its risk is detaching a proposition from the conditions under which it was observed or inferred. A statement such as ``the service accepts this argument'' is incomplete without a service version, scope, and time. A preference such as ``use the fast route'' may apply to one task and not to every future decision.

Separate factual assertions, user preferences, procedures, hypotheses, and authority. They have different update rules. A new authoritative correction can supersede a preference; conflicting observations may require uncertainty rather than replacement; an explicit prohibition cannot be overridden by repeated assistant-generated statements. A semantic store that deduplicates solely by embedding similarity can merge records whose surface meanings overlap while their roles differ. Deduplication should preserve distinctions needed for future interpretation.

A proposed semantic update represents competing propositions as a set with status rather than forcing immediate collapse to one string. For entity $x$, predicate $p$, and value $y$, retain the source, valid interval, and confidence of each assertion. A resolver produces a current view according to explicit rules. For example, source authority can precede recency, and recency can be evaluated within the correct valid-time interval. This is a semantic policy, not merely a database ordering operation.

Generative Agents and MemGPT show different ways of organizing external information and making selected state available to an underlying model~\cite{park2023generativeagents,packer2023memgpt}. Their architectures motivate the separation between stored experience and active context. They do not establish that extracted statements are universally correct, that biological memory mechanisms have been reproduced, or that an agent can retain all important information indefinitely. Those are stronger claims requiring different evidence.

Semantic extraction can also amplify certainty. A source says that a method ``may help'' under a particular setting; a first summary says it ``helps''; a later consolidated memory says it is the preferred method. Each individual rewrite can look reasonable while the chain removes uncertainty and scope. Provenance should therefore be attached to assertions, and the system should preserve the original qualification where it affects decisions. Repeated generation is not independent corroboration.

An evaluation can construct paired records differing only in modality, scope, or authority: an observation versus a hypothesis, a project preference versus a global preference, and a user correction versus a retrieved recommendation. The memory system should preserve the distinction through extraction, merging, retrieval, and use. This proposed test directly measures whether semantic normalization erases operational meaning. Exact-match fact recall alone does not expose these role errors.

